
\documentclass[11pt,a4paper]{article}

\usepackage[margin=1in]{geometry}
\usepackage{amsmath,amssymb}
\usepackage{siunitx}
\usepackage{booktabs}
\usepackage{array}
\usepackage{enumitem}
\usepackage{hyperref}
\usepackage{fancyhdr}
\usepackage{listings}
\usepackage{graphicx}

\hypersetup{
    colorlinks=true,
    linkcolor=blue,
    citecolor=blue,
    urlcolor=blue
}

\setlength{\headheight}{14pt}
\pagestyle{fancy}
\fancyhf{}
\lhead{Radioactivity --- AF December 2019}
\rhead{Radioactivity}
\cfoot{\thepage}

\title{\textbf{Radioactivity}\\Worked Solutions}
\author{AF --- December 2019}
\date{}

\begin{document}

\maketitle

\section*{Introduction}

This document gives worked solutions to the supplied BPhO
\emph{Radioactivity} problem sheet. The notation, constants, approximations and
methods follow the sheet and its accompanying solution document. Numerical
values have been evaluated explicitly so that the calculations can also be
reproduced computationally.

The central model is the radioactive decay law
\[
\frac{dN}{dt}=-\lambda N,
\]
with solution
\[
N(t)=N_0e^{-\lambda t}.
\]
The activity is
\[
A(t)=-\frac{dN}{dt}=\lambda N(t),
\]
and the half-life satisfies
\[
\boxed{\lambda=\frac{\ln 2}{t_{1/2}}}.
\]
Consequently,
\[
\boxed{N(t)=N_0\,2^{-t/t_{1/2}}}
\qquad\text{and}\qquad
\boxed{A(t)=A_0\,2^{-t/t_{1/2}}}.
\]

% ============================================================
\section{Question 1}

\subsection{(i) Half-life of radon}

The radon sample falls to $2\%$ of its original amount after $22$ days:
\[
\frac{N}{N_0}=0.02.
\]

Using
\[
\frac{N}{N_0}=2^{-t/t_{1/2}},
\]
we obtain
\[
0.02=2^{-22/t_{1/2}}.
\]

Taking logarithms,
\[
\frac{22}{t_{1/2}}
=
\log_2(50),
\]
so
\[
t_{1/2}
=
\frac{22}{\log_2(50)}
=
\boxed{3.90\ \mathrm{days}}.
\]

\subsection{(ii) Carbon-14 dating}

For fresh biomatter,
\[
A_0=238\ \mathrm{Bq\,kg^{-1}},
\]
whereas the hammer has
\[
A=1.05\ \mathrm{Bq\,kg^{-1}}.
\]

Since activity follows the same exponential law as the number of nuclei,
\[
\frac{A}{A_0}
=
2^{-t/t_{1/2}}.
\]

Therefore
\[
\frac{1.05}{238}
=
2^{-t/5370}.
\]

Hence
\[
t
=
5370
\log_2\left(\frac{238}{1.05}\right)
\approx
\boxed{4.20\times10^4\ \mathrm{years}}.
\]

Thus the hammer dates to approximately
\[
\boxed{42\,000\ \mathrm{years}}.
\]

\subsection{(iii) Attenuation of Sr-90 beta radiation}

The intensity attenuation law is
\[
I=I_0e^{-\mu x}.
\]

Here
\[
I_0=200\ \mathrm{Bq},
\qquad
I=10\ \mathrm{Bq},
\qquad
x=2.00\ \mathrm{mm}=2.00\times10^{-3}\ \mathrm m.
\]

Therefore
\[
\frac{10}{200}
=
e^{-\mu(0.002)}.
\]

Thus
\[
\mu
=
\frac{\ln(20)}{0.002}
\approx
\boxed{1.50\times10^3\ \mathrm{m^{-1}}}.
\]

The half-thickness is defined by
\[
\frac{I}{I_0}=\frac12,
\]
so
\[
\frac12=e^{-\mu x_{1/2}},
\]
giving
\[
\boxed{x_{1/2}=\frac{\ln2}{\mu}}
\]
and hence
\[
\boxed{x_{1/2}\approx4.63\times10^{-4}\ \mathrm m
=0.463\ \mathrm{mm}}.
\]

\subsection{(iv) Mass of Ra-226}

The activity is
\[
A=\lambda N,
\]
so
\[
N=\frac{A}{\lambda}
=
\frac{At_{1/2}}{\ln2}.
\]

Using
\[
A=9900\ \mathrm{Bq},
\qquad
t_{1/2}=1602\ \mathrm{yr}
\approx1602(365.25)(24)(3600)\ \mathrm s,
\]
gives
\[
N\approx7.22\times10^{14}\ \text{atoms}.
\]

The mass of one Ra-226 atom is approximately
\[
m_{\mathrm{Ra}}
=
226u,
\qquad
u=1.661\times10^{-27}\ \mathrm{kg}.
\]

Hence
\[
m
=
N(226u)
\approx
\boxed{2.71\times10^{-10}\ \mathrm{kg}}
\]
or
\[
\boxed{0.271\ \mathrm{mg}}.
\]

For U-235,
\[
t_{1/2}=2.22\times10^{16}\ \mathrm s,
\]
so
\[
N_U
=
\frac{A t_{1/2}}{\ln2}
\approx
3.17\times10^{20}.
\]

Therefore
\[
m_U
=
N_U(235u)
\approx
\boxed{1.24\times10^{-4}\ \mathrm{kg}}
=
\boxed{124\ \mathrm{mg}}.
\]

The much longer half-life of U-235 means that many more atoms are required to
produce the same activity.

\subsection{(v) Age of the Pa-231 statuette}

The Pa-231 initially decays into Ac-227, whose subsequent decay chain eventually
produces stable Pb-207. Under the problem's approximation, the measured fraction
of remaining Pa-231 is therefore
\[
\frac{N}{N_0}=0.10.
\]

Using
\[
\frac{N}{N_0}=2^{-t/t_{1/2}},
\]
with
\[
t_{1/2}=32800\ \mathrm{yr},
\]
gives
\[
0.1=2^{-t/32800}.
\]

Thus
\[
t
=
32800\log_2(10)
\approx
\boxed{1.09\times10^5\ \mathrm{years}}.
\]

Hence the statuette has been in space for approximately
\[
\boxed{109\,000\ \mathrm{years}}.
\]

% ============================================================
\section{Question 2 --- Iodine-123}

\subsection{(i) Decay equation}

Iodine has atomic number $Z=53$. Electron capture reduces the atomic number by
one without changing the mass number:
\[
{}^{123}_{53}\mathrm{I}+e^-
\rightarrow
{}^{123}_{52}\mathrm{Te}+\nu_e.
\]

The daughter nucleus is Te-123.

\subsection{(ii) Medical usefulness}

I-123 has several properties that make it useful for thyroid imaging:

\begin{itemize}
    \item its half-life of $13.22$ hours is long enough for preparation and imaging,
    but short enough that it does not remain radioactive for an excessively long
    period;
    \item it decays by electron capture and produces a useful $\gamma$ ray;
    \item the emitted $159\ \mathrm{keV}$ gamma radiation can escape the body and
    be detected externally;
    \item iodine is naturally concentrated by the thyroid, so the radioactive
    tracer preferentially reaches the tissue being imaged.
\end{itemize}

The radiation therefore provides a signal while allowing the distribution of
iodine in the thyroid to be measured.

\subsection{(iii) Activity after 24 hours}

The initial activity is
\[
A_0=2.5\times10^7\ \mathrm{Bq}.
\]

Using
\[
A(t)=A_0\,2^{-t/t_{1/2}},
\]
with
\[
t=24\ \mathrm h,
\qquad
t_{1/2}=13.22\ \mathrm h,
\]
we obtain
\[
A(24)
=
2.5\times10^7
2^{-24/13.22}.
\]

Therefore
\[
\boxed{A(24)\approx7.10\times10^6\ \mathrm{Bq}}.
\]

The activity actually measured in the thyroid will be lower because not all of
the administered iodine is taken up and retained by the thyroid. Some activity
remains elsewhere in the body and some is removed biologically.

\subsection{(iv) Number of I-123 atoms}

The decay constant is
\[
\lambda
=
\frac{\ln2}{13.22(3600)}
\approx1.456\times10^{-5}\ \mathrm{s^{-1}}.
\]

Since
\[
A_0=\lambda N_0,
\]
\[
N_0=\frac{A_0}{\lambda}.
\]

Hence
\[
N_0
=
\frac{2.5\times10^7}{1.456\times10^{-5}}
\approx
\boxed{1.72\times10^{12}\ \text{atoms}}.
\]

\subsection{(v) Total energy from gamma emission}

Each decay produces a gamma photon of energy
\[
E_\gamma=159\ \mathrm{keV}.
\]

Converting to joules,
\[
E_\gamma
=
159\times10^3(1.602\times10^{-19})
\approx
2.55\times10^{-14}\ \mathrm J.
\]

Assuming one such gamma emission per decay and complete absorption,
\[
E_{\rm total}=N_0E_\gamma.
\]

Thus
\[
E_{\rm total}
\approx
(1.72\times10^{12})(2.55\times10^{-14})
\approx
\boxed{4.37\times10^{-2}\ \mathrm J}.
\]

This is the total gamma energy associated with the initial population.

\subsection{(vi) Initial and average gamma power}

The instantaneous power immediately after administration is
\[
P_0=A_0E_\gamma.
\]

Thus
\[
P_0
=
(2.5\times10^7)
(159\times10^3)(1.602\times10^{-19})
\]
and therefore
\[
\boxed{P_0\approx6.37\times10^{-7}\ \mathrm{W}}.
\]

Over $48$ hours, the exact total emitted energy is
\[
E_{48}
=
N_0E_\gamma
\left(1-2^{-48/13.22}\right).
\]

Hence the average power is
\[
P_{\rm avg}
=
\frac{E_{48}}{48(3600)}
\approx
\boxed{2.33\times10^{-7}\ \mathrm W}.
\]

Thus the initial gamma-emission power is approximately
\[
\boxed{2.7P_{\rm avg}}.
\]

% ============================================================
\section{Question 3 --- Americium-241}

\subsection{(i) Decay equation}

Americium has atomic number $95$. Alpha decay reduces the mass number by $4$
and the atomic number by $2$:
\[
\boxed{
{}^{241}_{95}\mathrm{Am}
\rightarrow
{}^{237}_{93}\mathrm{Np}
+
{}^{4}_{2}\mathrm{He}
}.
\]

\subsection{(ii) Activity of a fresh smoke detector}

The problem states that the active material is $0.33\,\mu\mathrm g$ of
AmO$_2$. Using the supplied oxygen mass number $16$, the molar mass in the
problem's approximation is
\[
M_{\mathrm{AmO_2}}=241+2(16)=273\ \mathrm{g\,mol^{-1}}.
\]

Thus
\[
n=\frac{0.33\times10^{-6}\ \mathrm g}{273\ \mathrm{g\,mol^{-1}}}
\approx1.21\times10^{-9}\ \mathrm{mol},
\]
and the number of Am-241 nuclei is
\[
N_0=nN_A
\approx
\boxed{7.29\times10^{14}}.
\]

The decay constant is
\[
\lambda
=
\frac{\ln2}{432.2(365.25)(24)(3600)}
\approx5.08\times10^{-11}\ \mathrm{s^{-1}}.
\]

Therefore
\[
A_0=\lambda N_0
\approx
\boxed{3.70\times10^4\ \mathrm{Bq}}.
\]

Using
\[
1\ \mathrm{Ci}=3.7\times10^{10}\ \mathrm{Bq},
\]
gives
\[
\boxed{A_0\approx1.00\times10^{-6}\ \mathrm{Ci}}.
\]

This uses the stated AmO$_2$ mass directly. Treating the entire
$0.33\,\mu\mathrm g$ as Am-241, as in the accompanying handwritten
calculation, gives about $4.2\times10^4$ Bq; the AmO$_2$ calculation is the
one consistent with the wording of the question.

\subsection{(iii) Ionizations and detector current}

Each alpha particle has energy
\[
E_\alpha=5.486\ \mathrm{MeV}.
\]

The ionization energy is
\[
E_i=34\ \mathrm{eV}.
\]

Thus the approximate number of ion pairs produced per alpha particle is
\[
n_i
=
\frac{5.486\times10^6}{34}
\approx
1.61\times10^5.
\]

Using the activity found in part (ii), the ion-pair production rate is
\[
\dot N_i=A_0n_i
\approx
(3.70\times10^4)(1.61\times10^5)
\approx
\boxed{5.97\times10^9\ \mathrm{s^{-1}}}.
\]

An ion pair produces both an electron and a positive ion. If both charge
carriers are collected, their contributions to the detector current add, so
\[
I=2e\dot N_i.
\]

Hence
\[
I
\approx
2(1.602\times10^{-19})(5.97\times10^9)
\approx
\boxed{1.91\times10^{-9}\ \mathrm A}.
\]

\subsection{(iv) Estimating the molar mass of smoke particles}

Assume thermal equilibrium. The mean translational kinetic energy of a gas
particle is
\[
\frac12M\overline{v^2}
=
\frac32k_BT.
\]

Thus
\[
M\overline{v^2}=\text{constant}
\]
at fixed temperature, so
\[
M\propto\frac1{\overline{v^2}}.
\]

The current is proportional to the rate at which the charged particles are
collected. Following the model in the supplied solution, take the current
reduction to correspond to a reduction of the characteristic particle speed:
\[
\frac{v_{\rm smoke}}{v_{\rm air}}\approx0.2.
\]

Therefore
\[
\frac{\overline{v_{\rm smoke}^2}}
{\overline{v_{\rm air}^2}}
\approx0.2^2=0.04.
\]

Hence
\[
\frac{M_{\rm smoke}}{M_{\rm air}}
=
\frac{\overline{v_{\rm air}^2}}
{\overline{v_{\rm smoke}^2}}
\approx\frac1{0.04}=25.
\]

With
\[
M_{\rm air}=28.97\ \mathrm{g\,mol^{-1}},
\]
we obtain
\[
M_{\rm smoke}
\approx
25(28.97)
\approx
\boxed{7.24\times10^2\ \mathrm{g\,mol^{-1}}}.
\]

Thus the model predicts a characteristic smoke-particle molar mass of roughly
\[
\boxed{724\ \mathrm{g\,mol^{-1}}}.
\]

% ============================================================
\section{Question 4 --- The Elephant's Foot}

\subsection{(i) Absorbed particle rate}

The radiation dose rate is
\[
80\ \mathrm{Gy\,h^{-1}}
=
80\ \mathrm{J\,kg^{-1}\,h^{-1}}.
\]

A $10\ \mathrm{cm}\times10\ \mathrm{cm}\times10\ \mathrm{cm}$ cube of water has
volume
\[
V=(0.1)^3=10^{-3}\ \mathrm{m^3}.
\]

Taking the density of water to be approximately
\[
1000\ \mathrm{kg\,m^{-3}},
\]
the mass is
\[
m=1.0\ \mathrm{kg}.
\]

Therefore the absorbed energy per second is
\[
P
=
\frac{80\ \mathrm J}{3600\ \mathrm s}
=
2.22\times10^{-2}\ \mathrm{J\,s^{-1}}.
\]

For particle energy
\[
E_p=5\ \mathrm{MeV}
=
5\times10^6(1.602\times10^{-19})
\approx
8.01\times10^{-13}\ \mathrm J,
\]
the required particle rate is
\[
\dot N
=
\frac{P}{E_p}
\approx
\boxed{2.77\times10^{10}\ \mathrm{s^{-1}}}.
\]

The cube is at a distance from the source. Assuming approximately isotropic
spreading, the flux varies as
\[
\Phi\propto\frac1{r^2}.
\]

At three times the distance,
\[
\frac{\Phi'}{\Phi}
=
\frac1{3^2}
=
\boxed{\frac19}.
\]

Thus the particle flux is reduced by a factor of nine.

\subsection{(ii) Total activity of the Elephant's Foot}

The surface area of the hemispherical cap is
\[
A=\pi(a^2+h^2),
\]
with
\[
a=1.0\ \mathrm m,
\qquad
h=0.5\ \mathrm m.
\]

Hence
\[
A
=
\pi(1^2+0.5^2)
=
\boxed{3.93\ \mathrm{m^2}}.
\]

The particle rate found in part (i) is absorbed by an area
\[
A_{\rm cube}=(0.1)(0.1)=0.01\ \mathrm{m^2}.
\]

Therefore the estimated flux is
\[
\Phi
=
\frac{2.77\times10^{10}}{0.01}
=
2.77\times10^{12}\ \mathrm{m^{-2}s^{-1}}.
\]

Multiplying by the surface area,
\[
A_{\rm foot}
=
\Phi A
\]
gives
\[
A_{\rm foot}
\approx
(2.77\times10^{12})(3.93)
\approx
\boxed{1.09\times10^{13}\ \mathrm{Bq}}.
\]

This is an order-of-magnitude estimate based on the stated assumption that the
surface flux is uniform.

\subsection{(iii) Lead shielding}

For attenuation,
\[
I=I_0\,2^{-x/x_{1/2}},
\]
where
\[
x_{1/2}=4.2\ \mathrm{mm}.
\]

We require
\[
5
=
(2.77\times10^8)
2^{-x/4.2\mathrm{mm}}.
\]

Thus
\[
\frac{x}{4.2\ \mathrm{mm}}
=
\log_2\left(\frac{2.77\times10^8}{5}\right),
\]
so
\[
x
=
4.2
\log_2\left(\frac{2.77\times10^8}{5}\right)\ \mathrm{mm}.
\]

Therefore
\[
\boxed{x\approx108\ \mathrm{mm}}
\]
or approximately
\[
\boxed{10.8\ \mathrm{cm}}.
\]

% ============================================================
\section{Question 5 --- Pu-238 RTG}

\subsection{(i) Initial atoms, activity and thermal power}

The mass of plutonium is
\[
m=4.8\ \mathrm{kg}.
\]

Taking the molar mass of Pu-238 as approximately
\[
M=238\ \mathrm{g\,mol^{-1}},
\]
the number of moles is
\[
n=\frac{4800}{238}\approx20.17\ \mathrm{mol}.
\]

Thus
\[
N_0=nN_A
\approx
\boxed{1.21\times10^{25}\ \text{atoms}}.
\]

The half-life is
\[
t_{1/2}=87.7\ \mathrm{yr},
\]
so
\[
\lambda
=
\frac{\ln2}{87.7(365.25)(24)(3600)}
\approx
2.50\times10^{-10}\ \mathrm{s^{-1}}.
\]

The initial activity is therefore
\[
A_0=\lambda N_0
\approx
\boxed{3.04\times10^{15}\ \mathrm{Bq}}.
\]

Each alpha particle carries
\[
E_\alpha=5.593\ \mathrm{MeV}
\]
or
\[
E_\alpha
=
5.593\times10^6(1.602\times10^{-19})
\approx
8.96\times10^{-13}\ \mathrm J.
\]

Assuming every alpha particle's kinetic energy becomes thermal energy,
\[
P_0=A_0E_\alpha.
\]

Hence
\[
\boxed{P_0\approx2.73\times10^3\ \mathrm W}.
\]

The initial specific thermal power is
\[
\frac{P_0}{m}
=
\frac{2726}{4.8}
\approx
\boxed{568\ \mathrm{W\,kg^{-1}}}.
\]

\subsection{(ii) Heat power per kilogram as a function of time}

The activity obeys
\[
A(t)=A_0\,2^{-t/t_{1/2}}.
\]

Since the energy per decay is constant,
\[
P(t)=P_0\,2^{-t/t_{1/2}}.
\]

Therefore the power per kilogram is
\[
\boxed{
\frac{P(t)}{m}
=
568\,
2^{-t/87.7}
\ \mathrm{W\,kg^{-1}}
}.
\]

The graph is therefore an exponential decay with half-life $87.7$ years.

Useful reference points are
\[
\begin{array}{c|c}
t\ (\mathrm{yr}) & P/m\ (\mathrm{W\,kg^{-1}})\\
\hline
0 & 568\\
87.7 & 284\\
175.4 & 142\\
263.1 & 71.0\\
350.8 & 35.5
\end{array}
\]

The resulting exponential decay is plotted below.

\begin{center}
\includegraphics[width=0.82\textwidth]{pu238_power_per_kg.png}
\end{center}

The plotted curve uses the specific-power expression above and the reference
values in the table.

% ============================================================
\section{Question 6 --- Geiger--Nuttall law}

\subsection{(i) Constructing and testing the linear relation}

The Geiger--Nuttall law supplied in the sheet is
\[
\log_{10}\left(\frac{t_{1/2}}{\mathrm s}\right)
=
a+
\frac{bZ}{\sqrt{E_\alpha/\mathrm{MeV}}}.
\]

Define
\[
x=\frac{Z}{\sqrt{E_\alpha/\mathrm{MeV}}},
\qquad
y=\log_{10}\left(\frac{t_{1/2}}{\mathrm s}\right).
\]

The supplied spreadsheet calculates these quantities for the ten isotopes:

\begin{center}
\small
\begin{tabular}{lrrr}
\toprule
Isotope & $Z$ & $E_\alpha$ / MeV & $\log_{10}(t_{1/2}/\mathrm s)$\\
\midrule
Po-212 & 84 & 8.955 & $-6.52$\\
Po-214 & 84 & 7.680 & $-3.80$\\
Rn-218 & 86 & 7.263 & $-1.46$\\
At-217 & 85 & 7.000 & $-1.49$\\
Bi-213 & 83 & 5.870 & $3.44$\\
Ra-224 & 88 & 5.789 & $5.50$\\
Th-238 & 90 & 5.520 & $7.78$\\
Pa-231 & 91 & 5.150 & $12.01$\\
Np-237 & 93 & 4.959 & $13.83$\\
Th-232 & 90 & 4.081 & $17.65$\\
\bottomrule
\end{tabular}
\end{center}

A linear regression of $y$ against $x$ gives
\[
\boxed{
y=1.5077x-49.277
}
\]
with
\[
\boxed{R^2=0.9945}.
\]

The corresponding spreadsheet-style scatter plot and fitted line are shown
below.

\begin{center}
\includegraphics[width=0.82\textwidth]{geiger_nuttall.png}
\end{center}

The very high value of $R^2$ demonstrates that the supplied data strongly obey
the linearised Geiger--Nuttall relationship.

Thus
\[
\boxed{
\log_{10}\left(\frac{t_{1/2}}{\mathrm s}\right)
\approx
\frac{1.51Z}{\sqrt{E_\alpha/\mathrm{MeV}}}
-49.3
}.
\]

\subsection{(ii) Predicting the half-life of Francium-221}

For Francium-221,
\[
Z=87,
\qquad
E_\alpha=6.3\ \mathrm{MeV}.
\]

Therefore
\[
\log_{10}\left(\frac{t_{1/2}}{\mathrm s}\right)
\approx
\frac{1.51(87)}{\sqrt{6.3}}-49.3.
\]

Numerically,
\[
\frac{1.51(87)}{\sqrt{6.3}}
\approx52.34,
\]
so
\[
\log_{10}\left(\frac{t_{1/2}}{\mathrm s}\right)
\approx3.04.
\]

Hence
\[
t_{1/2}
\approx10^{3.04}\ \mathrm s
\approx
\boxed{1.09\times10^3\ \mathrm s}.
\]

In minutes,
\[
\boxed{t_{1/2}\approx18.2\ \mathrm{min}},
\]
or approximately
\[
\boxed{3.47\times10^{-5}\ \mathrm{yr}}.
\]

\section*{Summary}

The sheet develops a consistent computational picture of radioactive processes:

\[
\boxed{
\frac{dN}{dt}=-\lambda N
}
\qquad\Longrightarrow\qquad
\boxed{
N(t)=N_0e^{-\lambda t}
}
\]

with
\[
\boxed{
\lambda=\frac{\ln2}{t_{1/2}}
}.
\]

The same exponential law governs the number of nuclei, activity and the power
released by a radioactive source. Attenuation follows an analogous exponential
law,
\[
I=I_0e^{-\mu x},
\]
while the Geiger--Nuttall law reveals an approximately linear relationship between
$\log_{10}t_{1/2}$ and $Z/\sqrt{E_\alpha}$ for alpha emitters.

These relationships provide the basis for the numerical estimates throughout the
sheet, from radiocarbon dating and medical tracers to RTG power output, radiation
shielding and nuclear decay timescales.

\end{document}
